\documentclass{beamer}

\mode<presentation>
{
  \usetheme{Warsaw}
  \setbeamercovered{transparent}
}

\usetheme{Warsaw}
\setbeamersize{text margin left=4mm}
\setbeamersize{text margin right=4mm}

\usepackage[utf8]{inputenc}
\usepackage[backend=bibtex,style=authoryear]{biblatex}
\addbibresource{thesis.bib}

\usepackage{graphicx}
\usepackage{booktabs}
\usepackage{adjustbox}
\usepackage{capt-of}
\usepackage{threeparttable}
\usepackage{amsmath}
\usepackage{amssymb}
\usepackage{xcolor}
\usepackage{tikz}
\usepackage{appendixnumberbeamer}
\usepackage{dirtytalk}

\title[Dissertation Defense]{Essays on Heterogeneous Returns, Trust, and Household Wealth}
\author[Edwards]{Decory Edwards}
\institute[JHU]{Johns Hopkins University}
\date[June 22, 2026]{Final Dissertation Defense \\ June 22, 2026}

\newcommand{\assetpath}[1]{assets/#1}
\newcommand{\floatingnavbutton}[2]{%
  \begin{tikzpicture}[remember picture,overlay]
    \node[anchor=south east] at ([xshift=-0.5em,yshift=0.8em]current page.south east) {\hyperlink{#1}{\beamerbutton{#2}}};
  \end{tikzpicture}%
}
\newcommand{\centernavbutton}[2]{%
  \begin{center}
    \hyperlink{#1}{\beamerbutton{#2}}
  \end{center}
}
\newcommand{\slidetable}[1]{%
  \begin{center}
  \begin{adjustbox}{max totalsize={\textwidth}{0.70\textheight}}
  #1
  \end{adjustbox}
  \end{center}
}


\begin{document}

\begin{frame}
  \titlepage
\end{frame}

\begin{frame}{Roadmap}
  \tableofcontents[hideallsubsections]
\end{frame}

\section{Introduction}

\begin{frame}{Question and Unifying Theme}
  \begin{itemize}
    \item This dissertation asks how heterogeneous rates of return reshape the explanation, performance, and measurement of wealth inequality.
    \item The common theme is that wealth inequality is not only about who saves more; it is also about who earns different returns, why, and how those differences should be evaluated.
    \item Chapter 1 is structural macro, Chapter 2 is empirical household finance, and Chapter 3 is an alternative inequality measurement.
  \end{itemize}

  \vspace{0.5em}
  \textit{Core message: rates of return matter for understanding differences in household wealth.}
\end{frame}

\begin{frame}{How the Three Chapters Fit Together}
  \begin{enumerate}
    \item \textbf{Chapter 1}: Can return heterogeneity help a standard heterogeneous-agent model match empirical wealth moments?
    \item \textbf{Chapter 2}: Is there a systematic relationship between trust and returns to net wealth, and can we still estimate a persistent household return component?
    \item \textbf{Chapter 3}: If wealth is tied to social institutions and historical group structure, does the standard objective benchmark understate wealth inequality?
  \end{enumerate}

  \vspace{0.6em}
  \textit{Taken together, the chapters move from explaining wealth concentration, to explaining realized performance, to measuring inequality itself.}
\end{frame}

\section{Chapter 1}

\begin{frame}{Chapter 1: Matching Wealth Moments with Heterogeneous Returns}
  \begin{itemize}
    \item I incorporate return heterogeneity into a standard heterogeneous-agent model and estimate the return distribution needed to match the wealth distribution.
    \item Labor income uncertainty and life-cycle structure explain part of wealth dispersion, but not enough to reproduce the empirical skewness of wealth.
    \item Adding return heterogeneity yields a much better fit to SCF wealth moments and implies a more realistic aggregate MPC.
  \end{itemize}

  \vspace{0.6em}
  Empirically, this is motivated by evidence that realized household returns vary widely across people \parencite{aflgdmlp20, Daminato2024}.
\end{frame}

\begin{frame}[label=ch1-results]{Chapter 1: Main Result}
  \begin{columns}
    \column{0.56\textwidth}
    \begin{center}
      \includegraphics[width=\textwidth,height=0.65\textheight,keepaspectratio,trim={0.8in 3.4in 0.8in 5.0in},clip]{\assetpath{ch1_fig1_page31.pdf}}
    \end{center}
    \scriptsize\textit{Heterogeneous returns produce a Lorenz curve much closer to the SCF target.}

    \column{0.44\textwidth}
    \scriptsize
    \begin{itemize}
      \item Estimated uniform return distribution:
      \[
      R = 1.0204, \qquad \nabla = 0.06833.
      \]
      \item Top 5\% wealth share:
      \begin{itemize}
        \item Data: 57.4\%
        \item No heterogeneity: 23.3\%
        \item Heterogeneity: 58.8\%
      \end{itemize}
    \end{itemize}

  \end{columns}
\end{frame}

\begin{frame}[label=ch1-policy]{Chapter 1: Policy Implications — Wealth Distribution}
  \begin{columns}
    \column{0.55\textwidth}
    \begin{center}
      \includegraphics[width=\textwidth,height=0.68\textheight,keepaspectratio]{\assetpath{ch1_lorenz_tax_py.png}}
    \end{center}

    \column{0.45\textwidth}
    \scriptsize
    \begin{itemize}
      \item Two revenue-equivalent policies (each raising 1\% of GDP): wealth tax $1+r-\tau_w$ lowers the gross return for all; capital income tax $1+(1-\tau_{ci})r$ taxes only positive returns.
      \item In partial equilibrium, the CIT reduces wealth concentration modestly more than the WT in both the PY and LC calibrations.
      \item Intuition: low-return households face near-zero CIT (little positive capital income to tax) but a positive WT regardless of their return.
    \end{itemize}
  \end{columns}
\end{frame}

\begin{frame}[label=ch1-welfare]{Chapter 1: Policy Implications — Welfare}
  \scriptsize
  \begin{itemize}
    \item Welfare is measured by consumption-equivalent $\Delta$: the \% increase in consumption under the WT that makes a newborn --- before drawing their return type --- indifferent to the CIT. $\Delta > 0$ means CIT preferred.
    \item \textbf{Aggregate:} $\Delta = 0.20\%$ (PY) and $0.15\%$ (LC) --- the CIT is preferred on average for a newborn.
    \item \textbf{By return type:} 6 of 7 types prefer the CIT; only the highest-return type ($R=1.079$) prefers the WT. Low-return types gain most from switching to CIT because the WT reduces their already-low gross return while generating little revenue from them.
  \end{itemize}
  \vspace{0.3em}
  \slidetable{
  \begin{tabular}{lccccccc}
    \toprule
    & \textbf{Type 1} & \textbf{Type 2} & \textbf{Type 3} & \textbf{Type 4} & \textbf{Type 5} & \textbf{Type 6} & \textbf{Type 7} \\
    \midrule
    Baseline $R$ & 0.962 & 0.981 & 1.001 & 1.020 & 1.040 & 1.059 & 1.079 \\
    CE $\Delta$ (WT vs CIT) & +0.23\% & +0.27\% & +0.34\% & +0.32\% & +0.30\% & +0.24\% & $-$0.31\% \\
    \bottomrule
  \end{tabular}
  }
\end{frame}


\section{Chapter 2}

\begin{frame}{Chapter 2: Moderate Trust and Maximum Performance on Financial Investments}
  \begin{itemize}
    \item Economic performance is hump-shaped in trust: moderate trust maximizes outcomes at the individual level.
    \item Using the RAND HRS panel --- which includes both household financial data and a generalized trust measure --- we test whether this holds for returns to net wealth.
    \item It does: the return-maximizing trust level is 6--7 on a 1--10 scale, and the panel structure allows us to estimate a persistent household return component.
  \end{itemize}
\end{frame}

\begin{frame}[label=ch2-results]{Chapter 2: Main Result}
  \scriptsize
  \begin{itemize}
    \item Trust and Trust$^2$ are jointly significant: returns to net wealth are hump-shaped in trust.
    \item The result is robust across pooled OLS and correlated random effects (CRE/Mundlak) specifications.
  \end{itemize}
  \vspace{0.2em}
  \slidetable{
  \begin{tabular}{lcccc}
    \toprule
    & \multicolumn{2}{c}{Pooled OLS} & \multicolumn{2}{c}{CRE} \\
    \cmidrule(lr){2-3}\cmidrule(lr){4-5}
    & (1) Baseline & (2) Extended & (3) Baseline & (4) Extended \\
    \midrule
    Trust     & $0.03^{***}$ & $0.03^{***}$ & $0.03^{***}$ & $0.03^{***}$ \\
              & $(0.01)$     & $(0.01)$     & $(0.01)$     & $(0.01)$     \\
    Trust$^2$ & $-0.00^{**}$ & $-0.00^{**}$ & $-0.00^{***}$& $-0.00^{***}$\\
              & $(0.00)$     & $(0.00)$     & $(0.00)$     & $(0.00)$     \\
    \midrule
    $N$            & 2{,}919 & 2{,}899 & 2{,}919 & 2{,}899 \\
    $R^2$          & 0.04    & 0.04    & 0.15    & 0.15    \\
    $\rho$         &         &         & 0.16    & 0.15    \\
    $\sigma_u$     &         &         & 0.16    & 0.15    \\
    $\sigma_e$     &         &         & 0.36    & 0.36    \\
    \bottomrule
    \multicolumn{5}{l}{\textit{Controls: age, wealth, education, year; col.\ 2 \& 4 add demographics.}}
  \end{tabular}}
\end{frame}

\begin{frame}[label=ch2-trustmax]{Chapter 2: Return-Maximizing Trust Level}
  \begin{columns}
    \column{0.55\textwidth}
    \begin{center}
      \includegraphics[width=\textwidth,height=0.65\textheight,keepaspectratio]{\assetpath{ch2_trustdist_main.png}}
    \end{center}

    \column{0.45\textwidth}
    \scriptsize
    \slidetable{
    \begin{tabular}{lcc}
      \toprule
      Specification & No controls & With controls \\
      \midrule
      Avg.\ returns & 6.70 & 7.07 \\
      Pooled OLS    & 6.89 & 7.61 \\
      CRE           & 6.34 & 6.59 \\
      \bottomrule
    \end{tabular}}
    \vspace{0.4em}
    \begin{itemize}
      \item Estimated return-maximizing trust is consistently 6--7 across all specifications.
      \item The observed trust distribution concentrates near this optimal range.
    \end{itemize}
  \end{columns}
\end{frame}

\begin{frame}[label=ch2-persistent]{Chapter 2: Persistent Return Component}
  \scriptsize
  $\rho\approx0.16$ (intraclass correlation), $\sigma_u\approx0.16$ (persistent SD), $\sigma_e\approx0.36$ (idiosyncratic SD): roughly 16\% of total return variance reflects stable household differences, unchanged by adding portfolio-share controls.
  \vspace{0.3em}
  \begin{columns}
    \column{0.50\textwidth}
    \begin{center}
      \includegraphics[width=\textwidth,height=0.55\textheight,keepaspectratio]{\assetpath{ch2_fe_hist_main.png}}\\[2pt]
      \textit{FE: estimated individual effects (full controls).}
    \end{center}
    \column{0.50\textwidth}
    \begin{center}
      \includegraphics[width=\textwidth,height=0.55\textheight,keepaspectratio]{\assetpath{ch2_cre_hist_main.png}}\\[2pt]
      \textit{CRE: estimated individual effects (full controls).}
    \end{center}
  \end{columns}
\end{frame}

\section{Chapter 3}

\begin{frame}{Chapter 3: Wealth Inequality Under Ambiguity}
  \begin{itemize}
    \item Motivated by Rawls' discussion of justice behind a veil of ignorance, I propose a wealth inequality measure based on choice under ambiguity.
    \item The key idea is that wealth distributions may need to be ranked using socially relevant states, rather than objective uncertainty alone.
    \item I then adapt the social-welfare approach to inequality measurement to define an ambiguity-based wealth inequality measure.
  \end{itemize}

  \vspace{0.6em}
  \textit{Question: does the standard objective benchmark understate wealth inequality when racial structure is central?}
\end{frame}

\begin{frame}[label=ch3-setup]{Chapter 3: Inequality Measurement --- Setup}
  \scriptsize
  \textbf{Input:} A wealth frequency distribution $p(y)$ --- for each wealth level $y \in [0,\bar{y}]$, $p(y)$ is the share of the population at that wealth level. The objects being ranked are \textit{full distributions}, not individual outcomes.

  \vspace{0.5em}
  \textbf{Objective (Atkinson) measure.} A social planner with welfare function $U(y)$ (increasing, concave) ranks distributions by expected utility:
  $$W^O = \int_0^{\bar{y}} U(y)\, p(y)\, dy.$$
  The \textit{equally distributed equivalent} wealth $w_{EDE}$ satisfies $U(w_{EDE}) = W^O$. The inequality index is:
  $$I^O = 1 - \frac{w_{EDE}}{\mu}, \qquad \mu = \text{mean wealth.}$$

  \vspace{0.2em}
  Using the CRRA welfare function $U(y) = \frac{y^{1-\epsilon}}{1-\epsilon}$ (with $\epsilon \geq 0$ governing inequality aversion) yields a mean-independent, unit-free measure $I^O \in [0,1]$.
\end{frame}

\begin{frame}[label=ch3-ambiguity-derive]{Chapter 3: The Ambiguity-Based Measure}
  \scriptsize
  \textbf{Extension.} States $S = \{\text{Black},\,\text{White}\}$. The planner does not know the objective probability over states --- this is \textit{ambiguity} (Gilboa--Schmeidler 1989). Each act $f \in F$ maps a state to a wealth distribution: $f(s) = p(s)(y)$.

  \vspace{0.4em}
  The ambiguity-averse welfare criterion takes the worst-case weighting over beliefs $\lambda \in \Delta(S)$:
  $$W^A = \min_{\lambda \in \Delta(S)} \int_S \int_0^{\bar{y}} U(y)\,p(s)(y)\,dy\,d\lambda \;=\; \min_{s\,\in\, S}\;\mathbb{E}_{p(s)}\bigl[U(y)\bigr].$$
  The minimum puts all weight on whichever racial group has the lower within-group expected utility that year.

  \vspace{0.4em}
  The ambiguity EDE solves $U(w^A_{EDE}) = W^A$, giving:
  $$I^A = 1 - \frac{w^A_{EDE}}{\mu}.$$

  \vspace{0.2em}
  Since $W^A \leq W^O$ by construction, $I^A \geq I^O$: the ambiguity criterion \textbf{systematically reports higher inequality} than the objective benchmark for the same wealth distribution.
\end{frame}

\begin{frame}{Chapter 3: Data Pattern}
  \begin{columns}
    \column{0.70\textwidth}
    \begin{center}
      \includegraphics[width=\textwidth,height=0.72\textheight,keepaspectratio,trim={0.8in 5.75in 0.8in 2.15in},clip]{\assetpath{ch3_gap_page164.pdf}}
    \end{center}

    \column{0.30\textwidth}
    \scriptsize
    \begin{itemize}
      \item Persistent, large racial wealth gap in both net wealth and gross assets.
      \item This is the empirical benchmark the inequality measures speak to.
    \end{itemize}
  \end{columns}
\end{frame}

\begin{frame}[label=ch3-results]{Chapter 3: Ambiguity-Based Wealth Inequality}
  \begin{columns}
    \column{0.60\textwidth}
    \begin{center}
      \includegraphics[width=\textwidth,height=0.30\textheight,keepaspectratio]{\assetpath{ch3_measures_net.png}}\\[2pt]
      \includegraphics[width=\textwidth,height=0.30\textheight,keepaspectratio]{\assetpath{ch3_measures_gross.png}}
    \end{center}

    \column{0.40\textwidth}
    \scriptsize
    \begin{itemize}
      \item The objective benchmark understates wealth inequality relative to the ambiguity-based measure for the same wealth distribution.
      \item This holds for both net wealth (top) and gross assets (bottom), and the measures track the broad time pattern of the median wealth-gap diagnostics.
    \end{itemize}
  \end{columns}
\end{frame}

\section{Conclusion}

\begin{frame}{Conclusion}
  \begin{itemize}
    \item \textbf{Chapter 1}: return heterogeneity can match key wealth moments in a structural macro model, with an estimated distribution close to empirical evidence.
    \item \textbf{Chapter 2}: returns to net wealth exhibit a hump-shaped relationship with trust, and the maximizing level of trust lies in a moderate range between 6 and 7.
    \item \textbf{Chapter 3}: the ambiguity-based wealth inequality measure implies more inequality than the standard objective benchmark for the same wealth distribution.
  \end{itemize}

  \vspace{0.6em}
  \textit{Overall conclusion: heterogeneous returns are central both to explaining household wealth inequality and to evaluating it.}
\end{frame}

\end{document}
